\chapter{1951:Edwin M. McMillan and Glenn T. Seaborg}

\label{chap:chemistry1951}

The Nobel Prize in Chemistry for the year 1951 was awarded jointly to Edwin M. McMillan and Glenn T. Seaborg.

\textbf{Motivation:}

for their discoveries in the chemistry of the transuranium elements

\section{Edwin M. McMillan}

\subsection*{Early Life and Education}

Edwin M. McMillan was born on 1907-09-18 in Redondo Beach, California, USA.

\subsection*{Career and Major Research}

McMillan studied at the California Institute of Technology, where he received his doctorate in 1932, and then joined the Radiation Laboratory of the University of California, Berkeley, working with Ernest O. Lawrence's cyclotron group. In 1940 he discovered the first transuranium element, neptunium, together with Philip Abelson, by bombarding uranium with neutrons and showing that the resulting 2.3-day activity belonged to a new element of atomic number 93. He also introduced the principle of phase stability, which underlies the design of the synchrotron and of modern particle accelerators.

\subsection*{Nobel Prize-winning Work}

The work for which Edwin M. McMillan was awarded the Nobel Prize in Chemistry in 1951 was described by the Nobel Committee as: "for their discoveries in the chemistry of the transuranium elements".

The discovery of neptunium showed that elements heavier than uranium could be produced and identified in the laboratory, opening the study of the transuranium elements and establishing that the actinide series could be extended beyond uranium.

\subsection*{Significance and Impact}

McMillan's work demonstrated that the periodic table could be extended beyond uranium through neutron capture and beta decay. His discovery, and the plutonium work that followed from it in Berkeley, founded the chemistry of the transuranium elements and created the materials used in the first atomic weapons.

\subsection*{Later Career and Honors}

During the Second World War McMillan worked on radar at the Massachusetts Institute of Technology and at the Los Alamos Laboratory. After the war he returned to Berkeley and served as director of the Radiation Laboratory. He died on 1991-09-07 in El Cerrito, California.

\subsection*{Legacy}

McMillan is remembered as a discoverer of neptunium and as the originator of the principle of phase stability, which made possible the synchrotron and the large accelerators of modern particle physics.

\subsection*{Modern Developments}

McMillan's discovery of the transuranium elements neptunium and plutonium extended the periodic table and made the nuclear era possible. Plutonium became the fuel of modern nuclear reactors and weapons, while research on transuranics continues at facilities such as Berkeley and in the synthesis of superheavy elements. McMillan's work also included early insights into the synchrotron principle used in particle accelerators worldwide.

\subsection*{Selected Publications}

\begin{itemize}

\item \emph{Radioactive Element 93}, Physical Review, 1940.

\end{itemize}

\clearpage

\section{Glenn T. Seaborg}

\subsection*{Early Life and Education}

Glenn T. Seaborg was born on 1912-04-19 in Ishpeming, Michigan, USA.

\subsection*{Career and Major Research}

Seaborg studied at the University of California, Los Angeles, and took his doctorate at the University of California, Berkeley, in 1937. He spent his career at Berkeley, where in February 1941 he and his colleagues Joseph W. Kennedy and Arthur C. Wahl produced and isolated plutonium. He went on to co-discover several further transuranium elements, including americium and curium, and in 1944 proposed the actinide concept, placing the transuranium elements in a separate series in the periodic table.

\subsection*{Nobel Prize-winning Work}

The work for which Glenn T. Seaborg was awarded the Nobel Prize in Chemistry in 1951 was described by the Nobel Committee as: "for their discoveries in the chemistry of the transuranium elements".

Seaborg's discovery of plutonium and his isolation of the transuranium elements established a new branch of chemistry. His actinide concept correctly predicted the chemical behavior and placement of the elements beyond actinium in the periodic table.

\subsection*{Significance and Impact}

The chemistry of the transuranium elements founded by Seaborg became essential to nuclear energy, since plutonium-239 proved to be a fissionable material usable in nuclear reactors and weapons. His actinide concept organized the heaviest elements into a coherent group and guided the search for still heavier synthetic elements.

\subsection*{Later Career and Honors}

Seaborg was a professor at the University of California, Berkeley, and served as director of the Lawrence Berkeley Laboratory. He was also chairman of the United States Atomic Energy Commission from 1961 to 1971. Element 106, discovered in his laboratory, was named seaborgium in his honor, making him the first person to have an element named after him while still living. He died on 1999-02-25 in Lafayette, California.

\subsection*{Legacy}

Seaborg co-discovered ten transuranium elements and transformed the far end of the periodic table, and his name is preserved in the element seaborgium. His work remains central to the chemistry of the actinides and to the technology of nuclear power.

\subsection*{Modern Developments}

Seaborg's actinide concept reorganized the periodic table and guided the discovery of ten transuranium elements, from plutonium to seaborgium. These elements are central to nuclear energy, radiopharmaceuticals, and national security science, and Seaborg's periodic table reformulation is now standard in textbooks. Modern heavy-element and superheavy-element chemistry builds directly on his framework.

\subsection*{Selected Publications}

\begin{itemize}

\item \emph{Radioactive Element 94 from Deuterons on Uranium}, Physical Review, 1946.

\end{itemize}

\clearpage

\section*{Chapter Summary}

The 1951 prize recognized the discovery of the transuranium elements. Edwin M. McMillan and Philip Abelson produced the first of them, neptunium, in 1940, and Glenn T. Seaborg and his colleagues produced and isolated plutonium in 1941, work that established the chemistry of the elements beyond uranium and led to the actinide concept.

